\subsection{基础枚举与小模拟}
\subsubsection{模拟的核心是状态时刻}
模拟题要求按规则执行，并不意味着随意堆条件。先定义变量表示“操作之前”还是“操作之后”，再明确更新顺序。若同一轮的所有变化依赖旧状态，应把旧状态和新状态分开。若操作依次发生，则后面的操作通常应看到前面操作的结果。

枚举题还需说明为什么不重不漏。例如枚举无序数对，使用 $1\le i<j\le n$，每对只出现一次；若使用 $i\ne j$，则每对通常出现两次。判断条件正确和枚举范围正确同样重要。

\subsubsection{例题 网格的一轮同步更新}
\textbf{题目。}给定一个 $n\times m$ 的0/1网格。一次更新后，每个格子的新值由其旧状态的上下左右四邻格决定：值为1的邻格至少有两个，则新值为1，否则为0。越界邻格不计入。只进行一轮更新，$1\le n,m\le500$。

\textbf{输入。}第一行 $n,m$，随后每行 $m$ 个0或1。\textbf{输出。}更新后的网格。
\begin{verbatim}
输入
2 3
1 0 1
0 1 0
输出
0 1 0
1 0 1
\end{verbatim}
\textbf{推演。}左上格的两个有效邻格都是0，所以变为0；上方中间格的三个邻格都是1，所以变为1。所有计数都使用原数组 a，把答案放进 b。若直接覆盖 a，后面的格子会读到已经变化的数据。

\textbf{正确性。}程序对每格恰好检查四个方向，只有合法坐标参与计数，因此统计值与题意一致；写入 b 不改变 a，满足同步更新。时间和空间均为 $O(nm)$。
\Code{02_grid.cpp}

\subsubsection{矩形相交的边界}\ojpractice{CSP202303A}{田地丈量：第29次，2023-03，第1题；把每块田与固定区域求交后累加}
两个与坐标轴平行的矩形分别为 $[x_1,x_2]\times[y_1,y_2]$ 和 $[u_1,u_2]\times[v_1,v_2]$，且各端点已经有序。相交宽度与高度为
\[
w=\max(0,\min(x_2,u_2)-\max(x_1,u_1)),\qquad
h=\max(0,\min(y_2,v_2)-\max(y_1,v_1)).
\]
面积为 $wh$。例如 $[0,3]\times[0,2]$ 与 $[2,5]\times[1,4]$ 的交集宽、高都为1，面积为1。仅接触边界时面积为0。若题目统计整数格点而非连续面积，端点规则不同，不能擅自加1或减1。

\subsubsection{日期模拟}\ojpractice{CSP201509B}{日期计算：第5次，2015-09，第2题；用月份天数表消耗年内日序号}
日期问题先确定历法。公历闰年的判定是：能被400整除，或者能被4整除但不能被100整除。月份长度用表保存，二月按闰年调整。把日期加一天时，先增加日；若超过当月天数，日改为1并增加月；若月变成13，再进入下一年。每步只有一个明确定义的状态变化。

\begin{lstlisting}
bool leap(int y) {
    return y % 400 == 0 || (y % 4 == 0 && y % 100 != 0);
}
int days(int y, int m) {
    int d[] = {0,31,28,31,30,31,30,31,31,30,31,30,31};
    return d[m] + (m == 2 && leap(y));
}
\end{lstlisting}

\subsubsection{自测与解答}
\textbf{练习1。}对只有一个格子的网格，原值为1，更新后是多少？\textbf{解答。}为0，因为没有有效邻格。它检查越界和最小规模。

\textbf{练习2。}2000年和2100年是否闰年？\textbf{解答。}2000能被400整除，是闰年；2100能被100整除但不能被400整除，不是闰年。因此两年的二月分别有29天和28天。

\textbf{练习3。}矩形 $[0,2]\times[0,2]$ 与 $[2,4]\times[0,2]$ 的相交面积是多少？\textbf{解答。}宽度为0，面积为0；边接触不产生正面积。

\textbf{练习4。}若网格更新两轮，是否可以一直只计算 b？\textbf{解答。}不可以。第一轮结束后应交换 a、b，使第二轮从第一轮结果出发；每轮仍必须完整覆盖 b 的全部位置。
